\section{The current regularized DC-SBM result}\label{sec:sparse-restored}
\label{restore:algorithm}

Fix sufficiently small $a_0,\alpha>0$ and sufficiently large $\beta>0$,
depending only on the fixed structural parameter bounds. Choose them so
that the population degrees lie strictly inside the propensity-clamping
interval, and the population floor is inactive. Put
\[
 \mathbf d=A\one,\quad \bar D=n^{-1}\one^\top\mathbf d,\quad
 \tau=a_0\bar D,\quad G_i=\max(D_i,\tau).
\]
The all-zero graph has a fixed output. Otherwise $G$ is positive, even
when some vertices are isolated. Retain the $K$ eigenpairs with largest
absolute eigenvalues of $S_G=G^{-1/2}AG^{-1/2}$ and the raw degrees.
There is no requirement that a Perron degree vector be one of these
eigenvectors after flooring. Compute
\[
 R_0=G^{1/2}U\Lambda U^\top G^{1/2}=AH_G,
 \qquad H_G=G^{-1/2}UU^\top G^{1/2},
\]
and put, with $r=\diag(R_1)$ understood as a column vector,
\begin{equation}\label{restore:reconstruction}
 R_1=R_0+(\mathbf d-R_0\one)\one^\top/n,\qquad
 R^h=R_1+\frac{r\one^\top}{n-1}-\frac n{n-1}\diag(r).
\end{equation}
Thus $R^h\one=\mathbf d$ and $\diag(R^h)=0$ exactly. Only $(\mathbf d,U,\Lambda)$ are used
after encoding; $G$ is recovered from $\mathbf d$.

Initialize by normalizing the rows of $U$ and using the deterministic
radius-$1/8$, deletion-radius-$1/4$ count procedure in
\cref{dcsbm:initializer}. A zero row is assigned any fixed unit direction
for this initialization, without stopping other vertices.

Define
\[
 T_i=\min\{\max(D_i,\alpha\bar D),\beta\bar D\},\qquad
 \widehat\theta_i=T_i/\overline T_{\widehat z_i},\qquad
 \overline T_a=\widehat n_a^{-1}\sum_{j:\widehat z_j=a}T_j.
\]
Fit the symmetric block matrix by the exact signed Bernoulli working
criterion \eqref{dcsbm:block-criterion}, with $R^h$ and these degree
estimates, under the same fixed candidate-probability cap. Classify each
row by exact Bernoulli target profiling, with the degree parameter
profiled separately for every candidate community.

At least one fit-and-LRT update is performed. The procedure may use any
finite number of these hard refits, including finite data-dependent
stopping. The invariant basin below covers every such endpoint
simultaneously. An undefined empty-group or parameter-fit operation may
have a fixed fallback; those operations are well-defined on the proved
aggregate event. No individual zero degree triggers a graph-wide fallback.

\begin{theorem}[Regularized DC-SBM at every diverging expected-degree
scale]\label{restore:dc-theorem}\label{thm:regularized}
Under the fixed structural assumptions \eqref{eq:mean-model}--%
\eqref{eq:structural}, suppose
\[
 d=np\longrightarrow\infty,\qquad np^2\le C_1<\infty.
\]
The regularized algorithm of \cref{restore:algorithm}, using the
DC-SBM profile LRT, satisfies
\[
 \E\ell(\widehat z,z)\le C R_{\mathrm{prof},*}(n)
\]
for all sufficiently large $n$. The constant may depend on the fixed
structural bounds, clamping constants and $C_1$, but not on $n$.
Under the finite-$n$ equal-exposure condition
\eqref{eq:equal-exposure}, the right side may be replaced by
$C R_{\theta,*}(n)$. In particular, the classical balanced planted
partition with fixed separated $q/p$ is included. There is no lower
bound proportional to $\log n$ in this statement.
\end{theorem}


The reference $R_{\mathrm{prof},*}$ uses the full incident edge row.
The second comparison is available only under the stated equal-exposure
condition. The impossibility result in \cref{thm:imposs} shows why a
known-target-degree oracle is not a valid unrestricted benchmark.
Only retained degrees and one eigensystem enter the algorithm;
deleted graphs below are proof devices.
